% ============================================================
%  CAPÍTULO 6 – Metodología
% ============================================================
\chapter{Metodología}

\section{Modelo de desarrollo seleccionado}

Para el desarrollo de FitLabs se ha adoptado una metodología \textbf{incremental e iterativa}, inspirada en los principios ágiles de \textbf{Scrum}, adaptada a un equipo reducido de dos personas (José Luis Sánchez González y Carlos Fraidias del Valle) sin roles formales de \textit{Scrum Master} o \textit{Product Owner} diferenciados.

La elección de esta metodología se justifica por los siguientes motivos:

\begin{itemize}
  \item \textbf{Requisitos variables}: A lo largo del proyecto, las prioridades y el alcance de algunas funcionalidades se ajustaron en función del tiempo disponible y de la retroalimentación recibida del tutor.
  \item \textbf{Valor entregable temprano}: Al dividir el trabajo en incrementos, fue posible mostrar versiones funcionales de la aplicación desde las primeras semanas, lo que facilitó la validación del diseño y la arquitectura.
  \item \textbf{Flexibilidad}: La naturaleza del ciclo formativo, con entregas parciales en el primer y segundo trimestre, se adapta perfectamente al modelo iterativo.
  \item \textbf{Trabajo en paralelo}: Al ser dos desarrolladores, fue necesario definir bien los módulos de responsabilidad de cada uno, lo que encaja con la división de trabajo típica de sprints.
\end{itemize}

\section{Fases del ciclo de vida}

El proyecto se estructuró en las siguientes fases:

\subsection{Fase 1 -- Planificación (octubre 2025)}

Definición del alcance del proyecto, elección de tecnologías, propuesta técnica entregada el 24 de octubre de 2025. Se fijaron los módulos principales (Autenticación, Resumen del Día, Clientes, Rutinas, Calendario, Mensajes) y se creó el repositorio en GitHub con la estructura de ramas.

\subsection{Fase 2 -- Análisis de Requisitos (octubre -- noviembre 2025)}

Identificación y documentación de los requisitos funcionales y no funcionales. Se elaboraron los primeros prototipos en Figma para validar la interfaz de usuario antes de comenzar el desarrollo.

\subsection{Fase 3 -- Diseño (noviembre 2025)}

Diseño de la arquitectura de la aplicación (estructura modular de pantallas en Flutter), diseño del esquema de base de datos en Supabase (tablas, relaciones, políticas RLS) y prototipado de alta fidelidad en Figma completado en esta fase.

\subsection{Fase 4 -- Implementación (noviembre 2025 -- marzo 2026)}

Desarrollo iterativo por módulos. Cada módulo fue desarrollado, integrado y probado antes de avanzar al siguiente:

\begin{enumerate}
  \item Sistema de autenticación (Login / Registro).
  \item Pantalla de Resumen del Día.
  \item Módulo de Mis Clientes y Detalle de Cliente.
  \item Módulo de Crear Rutina con integración de API de ejercicios.
  \item Módulo de Calendario.
  \item Módulo de Mensajes.
\end{enumerate}

\subsection{Fase 5 -- Pruebas y Validación (marzo -- abril 2026)}

Pruebas funcionales de cada módulo, pruebas de integración entre componentes y validación de los requisitos definidos en la Fase 2. Se utilizaron dispositivos físicos Android para las pruebas.

\subsection{Fase 6 -- Documentación y Entrega (abril -- mayo 2026)}

Redacción de la memoria final, ajuste del código base, preparación de la presentación y empaquetado del material para la entrega del 4 de mayo de 2026.

\section{Herramientas de gestión y control}

\begin{itemize}
  \item \textbf{GitHub}: Control de versiones con ramas separadas por desarrollador (\texttt{Luis}, \texttt{Carlos}) y rama principal \texttt{main} — ramas correspondientes a José Luis Sánchez González y Carlos Fraidias del Valle respectivamente.
  \item \textbf{Figma}: Diseño de prototipos de alta fidelidad y gestión del sistema de diseño visual.
  \item \textbf{VS Code}: Entorno de desarrollo integrado principal para Flutter/Dart.
  \item \textbf{Supabase Dashboard}: Gestión visual de la base de datos, políticas RLS y autenticación de usuarios.
  \item \textbf{GitHub Copilot}: Asistencia en la redacción de código y documentación.
\end{itemize}
